\section{Safety}
\label{sec:safety}

\subsection{Intended use}
The \ST{} is a portable, battery powered tester for the electrical inspection of
\ac{FS} vehicles. It may be used to
\begin{itemize}
    \item supply a \ac{DUT} with a DC voltage (\qtyrange{0}{60}{\volt}) or DC current (\qtyrange{0}{10}{\ampere}),
    \item measure low resistances in 4-wire connection and resistances in 2-wire connection,
    \item measure insulation resistance with test voltages of \qtylist{125;250;500;1000}{\volt},
    \item measure DC voltage and DC current.
\end{itemize}
Any other use is not intended. The device may only be used by persons who are
qualified for work on electrical systems and have read this manual.
\todo[inline]{Define the required qualification (e.g. FS ESO/electrical scrutineer, electrically instructed person)}

\subsection{General safety instructions}

\begin{danger}
    In the \disp{Isolation Test} mode the test terminals carry up to \qty{1000}{\volt} DC.
    Do not touch the test leads, the terminals or the \ac{DUT} while a test is running.
    The \ac{DUT} can stay charged after the test.
    \todo[inline]{Confirm: is the DUT actively discharged after the isolation test (Q18)? How long to wait?}
\end{danger}

\begin{warning}
    \begin{itemize}
        \item Only use the device with undamaged enclosure, terminals and test leads.
        \item Only use test leads with shrouded \qty{4}{\milli\meter} plugs rated for the voltage of the measurement.
        \item Do not exceed the measurement category and the maximum input voltage (\autoref{sec:specs}).
        \item Do not connect an external voltage to the terminals in the source modes, the
              resistance modes, \disp{Ampmeter} and \disp{Isolation Test}.
        \item Do not use the device near flammable gases or vapours.
    \end{itemize}
    \todo[inline]{Measurement category (requirement: CAT III 600 V) not yet verified / certified}
\end{warning}

\begin{caution}
    \begin{itemize}
        \item Only charge the device as described in \autoref{sec:charging}.
        \item The device contains a \ac{LiFePO4} battery. Do not open, short-circuit or heat it.
        \item Operating temperature: \tbc{\qtyrange{0}{50}{\celsius}}. Protection against rain: \tbc{IP54}.
    \end{itemize}
\end{caution}

\subsection{Protective functions}
The device contains several independent protective functions. They protect
the device; they do not replace the safety rules for working on electrical systems.
\begin{itemize}
    \item \ac{BMS}: protects the battery against over-/undervoltage, overcurrent and overtemperature.
    \item Overvoltage protection: if an excessive voltage is detected at the output,
          the output is short-circuited and the device is locked until it is switched off
          and on again (\autoref{sec:messages}). In case of an external source this may blow
          the internal fuse.
    \item Overcurrent protection of the \ac{HV} output.
    \item Voltage and current limits per mode, temperature monitoring.
    \item Automatic power-off after \qty{10}{\minute} without operation.
\end{itemize}
\todo[inline]{Emergency off: describe the hardware switch-off via ESC (how long to hold?)}
